\section{Población, muestra y muestreo}

\subsection{Población}

La población documental está constituida por los registros de la Encuesta Nacional Agropecuaria del INEI y las operaciones de Aduanet para la subpartida 0804400000 disponibles en el periodo de estudio. El documento técnico de base reporta aproximadamente 5\,607 observaciones por año en la fuente agropecuaria y 128\,222 registros aduaneros antes de la depuración. Estas cantidades son referencias del origen y no se asumirán como el número final de observaciones válidas.

La población usuaria está conformada por productores, asesores y técnicos vinculados con la producción de palta Hass en La Libertad que puedan utilizar información de precios, rentabilidad y ciclos para apoyar decisiones de campaña.

\subsection{Muestra}

En el componente documental se aplicará un criterio censal: se conservarán todos los registros que cumplan los criterios de inclusión y calidad. La muestra analítica definitiva se informará después del proceso ETL, junto con el número de registros excluidos y su causa.

Para la evaluación del instrumento y del prototipo se realizará una prueba piloto con un mínimo de 30 participantes, de acuerdo con la guía UNT-2026. Esta cantidad corresponde a una evaluación inicial de confiabilidad y usabilidad; no sustituye una muestra probabilística representativa de todos los productores de La Libertad.

\subsection{Muestreo}

El muestreo documental será no probabilístico por criterios y cobertura censal de los registros accesibles. Se incluirán operaciones con fecha válida, valor FOB positivo, peso neto positivo, destino identificable y correspondencia con la subpartida seleccionada. Se excluirán duplicados, operaciones anuladas, valores imposibles y registros sin información suficiente para calcular el precio unitario.

La selección de participantes será no probabilística e intencional, considerando experiencia o participación en la cadena de palta Hass, disponibilidad para completar las dos mediciones y consentimiento informado. Se procurará incluir perfiles con diferentes tamaños de unidad productiva para ampliar la diversidad de la prueba piloto.
